This project explored three common attacks that exploit weaknesses in the Transmission Control Protocol: TCP SYN flooding, TCP reset attacks, and TCP session hijacking. Although each attack targets a different aspect of TCP, they all demonstrate how assumptions made during the protocol's original design can be abused when additional security mechanisms are not in place. Through a series of Docker-based demonstrations, each attack was successfully reproduced in a controlled environment, allowing the underlying protocol behavior to be observed and analyzed. The experiments also highlighted the importance of TCP concepts such as the three-way handshake, the connection four-tuple, and sequence numbers in both establishing reliable communication and enabling certain attacks.

In addition to demonstrating the attacks themselves, this project examined practical mitigation techniques and their effectiveness. Mechanisms such as SYN cookies, sequence number randomization, RFC 5961 challenge ACKs, anti-spoofing measures, and protections against on-path attacks all help reduce the risk posed by these vulnerabilities. While modern operating systems and network devices implement many of these defenses, the demonstrations show that TCP security still depends heavily on proper network configuration and layered security practices. Overall, the project provided a practical understanding of how TCP attacks operate, why they succeed, and how they can be detected and mitigated in real-world networks.

